% GENERATED FILE. Edit canonical lessons, then run npm run generate:books.
% canonical-source-hash: fnv1a64:2cc7046a3f05f3f2
% canonical-lessons: GE-C398-landstrasse, GE-C398-ausfahrt, GE-C398-auffahrt, GE-C398-kreisverkehr, GE-C398-abbiegen

\chapter{Über Land}
\label{ch:ge-uber-land}
\hlchaptermodality{398}

\begin{chapteropening}
\textbf{By the end of this chapter:} \emph{I can talk about roads outside town: the country road, the exit, the slip road, the roundabout and turning off.}

Complete the last lesson of chapter 398: I can talk about roads outside town: the country road, the exit, the slip road, the roundabout and turning off.
\end{chapteropening}

\section[die Landstraße]{die Landstraße — the country road, the main road (outside towns)}

\label{lesson:GE-C398-landstrasse}



\begin{quote}
Before the new one: say the German for at least, then the German for during.
\end{quote}

\subsection*{You'll want to know: die Landstraße}
\textbf{die Landstraße} — \textquotedblleft{}the country road, the main road (outside towns)\textquotedblright{}.

\begin{quote}
Wir fahren lieber auf der Landstraße als auf der Autobahn.
\end{quote}

\begin{quote}
\emph{We prefer to drive on the country road rather than on the motorway.}
\end{quote}

Plural: \emph{die Landstraßen}. Slower than the \emph{Autobahn}, often with villages on the way.

\begin{cousinweb}
\emph{Landstraße} is a compound of \emph{Land} 'land, open country', inherited from Proto-Germanic \emph{landą} (EtymDB), and \emph{Straße} 'street', an early borrowing from Latin \emph{(via) strāta} 'paved road' (Lewis \& Short) that reached Middle High German as \emph{strāze} (EtymDB).
\end{cousinweb}

\begin{grammarlens}[title={What you've built}]
One more word for telling, explaining and describing at B1.
\end{grammarlens}

\subsection*{Your turn}
\begin{itemize}
\raggedright
  \item \emph{Say it:} \emph{die Landstraße}
  \item \emph{Say it:} \emph{die Landstraße}, once more
  \item \emph{Say it:} a sentence of your own with \emph{die Landstraße}
  \item \emph{Recall:} say the German for at least, then the German for during, then say \emph{die Landstraße} again
\end{itemize}

\begin{tcolorbox}[breakable,colback=teal!4,colframe=teal!35!black,title={Before you move on}]
What does die Landstraße mean? (\textquotedblleft{}The country road, the main road (outside towns)\textquotedblright{}.) Say it once more.
\end{tcolorbox}
\section[die Ausfahrt]{die Ausfahrt — the exit (from a motorway or a garage)}

\label{lesson:GE-C398-ausfahrt}



\begin{quote}
Before the new one: say the German for during, then the German for the country road.
\end{quote}

\subsection*{You'll want to know: die Ausfahrt}
\textbf{die Ausfahrt} — \textquotedblleft{}the exit (from a motorway or a garage)\textquotedblright{}.

\begin{quote}
Wir nehmen die nächste Ausfahrt und fahren in die Stadt.
\end{quote}

\begin{quote}
\emph{We take the next exit and drive into town.}
\end{quote}

Plural: \emph{die Ausfahrten}. \emph{Ausfahrt freihalten} — keep the exit clear, a sign on many garage doors.

\begin{cousinweb}
\emph{Ausfahrt} is built from \emph{aus} 'out' and \emph{Fahrt} 'journey, drive', which comes from the inherited verb \emph{fahren} 'to travel', Proto-Germanic \emph{faraną} (EtymDB).
\end{cousinweb}

\begin{grammarlens}[title={What you've built}]
One more word for telling, explaining and describing at B1.
\end{grammarlens}

\subsection*{Your turn}
\begin{itemize}
\raggedright
  \item \emph{Say it:} \emph{die Ausfahrt}
  \item \emph{Say it:} \emph{die Ausfahrt}, once more
  \item \emph{Say it:} a sentence of your own with \emph{die Ausfahrt}
  \item \emph{Recall:} say the German for during, then the German for the country road, then say \emph{die Ausfahrt} again
\end{itemize}

\begin{tcolorbox}[breakable,colback=teal!4,colframe=teal!35!black,title={Before you move on}]
What does die Ausfahrt mean? (\textquotedblleft{}The exit (from a motorway or a garage)\textquotedblright{}.) Say it once more.
\end{tcolorbox}
\section[die Auffahrt]{die Auffahrt — the slip road (onto a motorway); the driveway}

\label{lesson:GE-C398-auffahrt}



\begin{quote}
Before the new one: say the German for the country road, then the German for the exit.
\end{quote}

\subsection*{You'll want to know: die Auffahrt}
\textbf{die Auffahrt} — \textquotedblleft{}the slip road (onto a motorway); the driveway\textquotedblright{}.

\begin{quote}
An der Auffahrt zur Autobahn steht ein langer Stau.
\end{quote}

\begin{quote}
\emph{There is a long traffic jam at the slip road onto the motorway.}
\end{quote}

Plural: \emph{die Auffahrten}. Also the drive up to a big house.

\begin{cousinweb}
\emph{Auffahrt} is built from \emph{auf} 'up, onto' and \emph{Fahrt} 'journey, drive', from the inherited verb \emph{fahren} 'to travel' (EtymDB). EtymDB lists the word but gives no older form.
\end{cousinweb}

\begin{grammarlens}[title={What you've built}]
One more word for telling, explaining and describing at B1.
\end{grammarlens}

\subsection*{Your turn}
\begin{itemize}
\raggedright
  \item \emph{Say it:} \emph{die Auffahrt}
  \item \emph{Say it:} \emph{die Auffahrt}, once more
  \item \emph{Say it:} a sentence of your own with \emph{die Auffahrt}
  \item \emph{Recall:} say the German for the country road, then the German for the exit, then say \emph{die Auffahrt} again
\end{itemize}

\begin{tcolorbox}[breakable,colback=teal!4,colframe=teal!35!black,title={Before you move on}]
What does die Auffahrt mean? (\textquotedblleft{}The slip road (onto a motorway); the driveway\textquotedblright{}.) Say it once more.
\end{tcolorbox}
\section[der Kreisverkehr]{der Kreisverkehr — the roundabout (US: traffic circle)}

\label{lesson:GE-C398-kreisverkehr}



\begin{quote}
Before the new one: say the German for the exit, then the German for the slip road.
\end{quote}

\subsection*{You'll want to know: der Kreisverkehr}
\textbf{der Kreisverkehr} — \textquotedblleft{}the roundabout (US: traffic circle)\textquotedblright{}.

\begin{quote}
Im Kreisverkehr nimmt der Bus die zweite Ausfahrt.
\end{quote}

\begin{quote}
\emph{At the roundabout the bus takes the second exit.}
\end{quote}

Plural: \emph{die Kreisverkehre}. In everyday speech also \emph{der Kreisel}.

\begin{cousinweb}
\emph{Kreisverkehr} is a compound of \emph{Kreis} 'circle', inherited through Middle and Old High German \emph{kreiz} (EtymDB), and \emph{Verkehr} 'traffic', built on \emph{kehren} 'to turn'. EtymDB does not list the compound itself.
\end{cousinweb}

\begin{grammarlens}[title={What you've built}]
One more word for telling, explaining and describing at B1.
\end{grammarlens}

\subsection*{Your turn}
\begin{itemize}
\raggedright
  \item \emph{Say it:} \emph{der Kreisverkehr}
  \item \emph{Say it:} \emph{der Kreisverkehr}, once more
  \item \emph{Say it:} a sentence of your own with \emph{der Kreisverkehr}
  \item \emph{Recall:} say the German for the exit, then the German for the slip road, then say \emph{der Kreisverkehr} again
\end{itemize}

\begin{tcolorbox}[breakable,colback=teal!4,colframe=teal!35!black,title={Before you move on}]
What does der Kreisverkehr mean? (\textquotedblleft{}The roundabout (US: traffic circle)\textquotedblright{}.) Say it once more.
\end{tcolorbox}
\section[abbiegen]{abbiegen — to turn off (left or right)}

\label{lesson:GE-C398-abbiegen}



\begin{quote}
Before the new one: say the German for the slip road, then the German for the roundabout.
\end{quote}

\subsection*{You'll want to know: abbiegen}
\textbf{abbiegen} — \textquotedblleft{}to turn off (left or right)\textquotedblright{}.

\begin{quote}
Nach der Kirche biegt der Bus links ab.
\end{quote}

\begin{quote}
\emph{After the church the bus turns left.}
\end{quote}

Separable: \emph{er biegt ab}. Perfekt with \emph{sein}: \emph{wir sind links abgebogen}.

\begin{cousinweb}
\emph{abbiegen} is built from \emph{ab} 'off, away' and \emph{biegen} 'to bend', which is inherited through Middle and Old High German (\emph{biegen}, \emph{biogan}) from Proto-Germanic \emph{beuganą} (EtymDB). English \emph{bow} 'to bend' is a cousin.
\end{cousinweb}

\begin{grammarlens}[title={What you've built}]
One more word for telling, explaining and describing at B1.
\end{grammarlens}

\subsection*{Your turn}
\begin{itemize}
\raggedright
  \item \emph{Say it:} \emph{abbiegen}
  \item \emph{Say it:} \emph{abbiegen}, once more
  \item \emph{Say it:} a sentence of your own with \emph{abbiegen}
  \item \emph{Recall:} say the German for the slip road, then the German for the roundabout, then say \emph{abbiegen} again
\end{itemize}

\begin{tcolorbox}[breakable,colback=teal!4,colframe=teal!35!black,title={Before you move on}]
What does abbiegen mean? (\textquotedblleft{}To turn off (left or right)\textquotedblright{}.) Say it once more.
\end{tcolorbox}
